\documentclass[aspectratio=169,10pt]{beamer}
\usepackage[T1]{fontenc}
\usepackage[utf8]{inputenc}
\usepackage[spanish,es-noshorthands]{babel}
\usepackage{lmodern,booktabs,tabularx,array}
\usepackage{tikz}
\usetikzlibrary{positioning,arrows.meta}
\definecolor{navy}{HTML}{173750}
\definecolor{teal}{HTML}{007F86}
\definecolor{muted}{HTML}{586D7B}
\definecolor{soft}{HTML}{EDF5F6}
\definecolor{orange}{HTML}{A15A18}
\definecolor{lightgray}{HTML}{F4F7F9}
\setbeamertemplate{navigation symbols}{}
\setbeamertemplate{footline}{\leavevmode\hbox{\begin{beamercolorbox}[wd=.86\paperwidth,ht=3ex,dp=1.3ex,leftskip=.65cm]{normal text}\scriptsize\color{muted}RetenIA \quad | \quad Kevin Andrés Erazo Yepez \quad | \quad Business Pitch\end{beamercolorbox}\begin{beamercolorbox}[wd=.14\paperwidth,ht=3ex,dp=1.3ex,center]{normal text}\scriptsize\color{teal}\insertframenumber/5\end{beamercolorbox}}}
\setbeamercolor{normal text}{fg=navy,bg=white}
\setbeamercolor{frametitle}{fg=navy,bg=white}
\setbeamerfont{frametitle}{size=\Large,series=\bfseries}
\setbeamercolor{block title}{bg=soft,fg=teal}
\setbeamercolor{block body}{bg=lightgray,fg=navy}
\setbeamercolor{alerted text}{fg=orange}
\setbeamertemplate{itemize item}{\color{teal}\small$\blacktriangleright$}
\newcommand{\eyebrow}[1]{\vspace{-0.18cm}{\footnotesize\bfseries\color{teal}#1}\par\vspace{0.1cm}}
\newcommand{\bigstat}[3]{\begin{minipage}[t]{.29\linewidth}\colorbox{soft}{\parbox{\dimexpr\linewidth-2\fboxsep}{\centering\vspace{.17cm}{\LARGE\bfseries\color{teal}#1}\par\vspace{.09cm}{\scriptsize\color{navy}#2}\vspace{.17cm}}}\\[.12cm]{\tiny\color{muted}#3}\end{minipage}}
\begin{document}

\begin{frame}[t]{La oportunidad: actuar antes de perder clientes}
\eyebrow{01 / EL PROBLEMA}
\vspace{.16cm}
{\large Una operadora no puede gastar el mismo presupuesto de retención en todos los clientes.}\par
\vspace{.43cm}
\bigstat{7.043}{clientes en Telco}{Dataset de demostración}\hfill
\bigstat{$\approx$26,5\%}{abandono histórico}{Clase minoritaria}\hfill
\bigstat{19}{variables predictoras}{Perfil y servicios del cliente}
\vspace{.47cm}
\begin{block}{La pregunta de negocio}
¿A quién contactar primero cuando los recursos comerciales son limitados?
\end{block}
\vspace{.15cm}
{\small\color{muted}El reto no es solo clasificar: también evitar contactos innecesarios y pérdidas por bajas no detectadas.}
\end{frame}

\begin{frame}[t]{RetenIA: predicción, explicación y prioridad}
\eyebrow{02 / LA SOLUCIÓN}
\begin{columns}[T,onlytextwidth]
\column{.52\textwidth}
\begin{tikzpicture}[node distance=0.37cm, every node/.style={font=\small}, box/.style={draw=teal, rounded corners=3pt, fill=soft, text width=6.2cm, minimum height=.68cm, align=left, inner sep=6pt}]
\node[box] (a) {\textbf{1. Diagnosticar} \quad Explorar segmentos};
\node[box, below=of a] (b) {\textbf{2. Estimar} \quad P(Churn=1) individual};
\node[box, below=of b] (c) {\textbf{3. Explicar} \quad SHAP y LIME por cliente};
\node[box, below=of c] (d) {\textbf{4. Priorizar} \quad Ordenar por presupuesto};
\draw[-{Latex[length=2mm]},teal,thick] (a)--(b);
\draw[-{Latex[length=2mm]},teal,thick] (b)--(c);
\draw[-{Latex[length=2mm]},teal,thick] (c)--(d);
\end{tikzpicture}
\column{.43\textwidth}
\begin{block}{Para equipos de CRM}
\begin{itemize}
\item Panel interactivo de riesgo.
\item Modelos exportados y reproducibles.
\item Justificación de cada predicción.
\item Simulador de contactos de retención.
\end{itemize}
\end{block}
\vspace{.20cm}
{\small\color{muted}El sistema apoya decisiones humanas. No reentrena dentro de Streamlit ni garantiza evitar bajas.}
\end{columns}
\end{frame}

\begin{frame}[t]{Evidencia real: desempeño y selección transparente}
\eyebrow{03 / VALIDACIÓN Y LIMITACIONES}
\vspace{.12cm}
\begin{center}
\renewcommand{\arraystretch}{1.32}
\begin{tabular}{lccc}
\toprule
\textbf{Modelo (test reservado)} & \textbf{F1} & \textbf{Recall} & \textbf{ROC-AUC}\\
\midrule
Logistic Regression & \textbf{0.604} & 0.559 & 0.842\\
Random Forest base & 0.553 & 0.505 & 0.818\\
NSGA-II (Pareto) & 0.581 & 0.535 & 0.830\\
\bottomrule
\end{tabular}
\end{center}
\vspace{.18cm}
\begin{block}{Selección previa a la evaluación final}
Configuración operativa: \textbf{NSGA-II (Pareto)}. Se aplicó el criterio F1-CV y, cuando correspondía, el compromiso con latencia de NSGA-II. F1-CV = \textbf{0.5608}.
\end{block}
\vspace{.10cm}
{\small\color{orange}\textbf{Evidencia independiente:} Logistic Regression logró F1=0.604, superior al modelo operativo F1=0.581; no alteramos la selección tras test.}

\end{frame}

\begin{frame}[t]{Optimizar dos objetivos y explicar predicciones}
\eyebrow{04 / HYPERPARAMETER OPTIMIZATION, NSGA-II + XAI}
{\footnotesize\color{teal}Grid Search \quad | \quad Random Search \quad | \quad Hyperband (Optuna) \quad | \quad NSGA-II}\par
\vspace{.07cm}
\begin{columns}[T,onlytextwidth]
\column{.48\textwidth}
\begin{block}{Frente de Pareto}
\centering\small $\max_{\theta}\;F1_{CV}(\theta)$ \quad y \quad $\min_{\theta}\;T_{infer}(\theta)$
\vspace{.12cm}
\begin{itemize}
\item Soluciones no dominadas: no hay mejora simultánea gratuita.
\item Se acepta hasta 0,01 de pérdida de F1 frente al máximo de Pareto para buscar menor latencia.
\item La latencia depende del equipo y de la carga.
\end{itemize}
\end{block}
\column{.48\textwidth}
\begin{block}{Top variables SHAP (global)}
\begin{tabular}{@{}lr@{}}
\textbf{Variable} & \textbf{Importancia}\\
\midrule
Contract & 0.120\\tenure & 0.050\\OnlineSecurity & 0.045\\InternetService & 0.045\\
\end{tabular}
\end{block}
\vspace{.09cm}
{\small LIME explica dos clientes locales; la importancia del bosque compara reducciones de impureza.}
\end{columns}
\vspace{.10cm}
{\scriptsize\color{muted}Random Search y Hyperband (Optuna) empataron en F1-CV=0.5736; Random Search fue más rápido.}\par
{\small\color{orange}\textbf{Precaución:} importancia predictiva no prueba causalidad.}
\end{frame}

\begin{frame}[t]{Oferta: piloto controlado antes de contratar}
\eyebrow{05 / PROPUESTA Y CIERRE}
\vspace{.20cm}
\begin{columns}[T,onlytextwidth]
\column{.50\textwidth}
\begin{block}{Modelo comercial propuesto}
\begin{itemize}
\item Cliente: gerencia de CRM y retención.
\item Producto: panel de riesgo, XAI y prioridad.
\item Licencia \textbf{hipotética}: \textbf{USD 59/mes}.
\item Costos, margen y eficacia son supuestos editables.
\end{itemize}
\end{block}
\column{.45\textwidth}
\begin{block}{Prueba que sí aporta evidencia}
\begin{enumerate}
\item Piloto A/B con grupo de control.
\item Medir bajas evitadas de forma incremental.
\item Estimar retorno neto con costos reales.
\item Revisar privacidad, sesgos y consentimiento.
\end{enumerate}
\end{block}
\end{columns}
\vspace{.36cm}
\begin{center}
\colorbox{soft}{\parbox{.90\linewidth}{\centering\vspace{.20cm}\large\bfseries\color{teal}¿Podemos evaluar un piloto medible antes de escalar?\vspace{.20cm}}}
\end{center}
\vspace{.10cm}
{\tiny\color{muted}Cifras predictivas: validación independiente del modelo tras nueva ejecución. Simulación comercial: supuestos no observados en Telco.}
\end{frame}
\end{document}
